% !TEX root = ../main.tex
\section{Conclusions}\label{sec:conclusions}



In this work, we have considered a GeV-scale or lighter long-lived lightest neutralino, 
which is necessarily bino-like, in the minimal R-parity-violating (RPV) supersymmetric model.
We have focused on lepton-number-violating operators in the RPV superpotential: $LH_u$, 
$LL\bar{E}$, and $LQ\bar{D}$. Such light neutralinos are still allowed by all 
experimental and observational constraints, as long as they decay, for instance via RPV couplings, 
so as to avoid overclosing the Universe. Since these RPV couplings are bounded to be small, such 
light binos, which we assume to be the lightest supersymmetric particle in our theory, are expected 
to have a relatively long lifetime. Via the considered couplings, the binos can decay leptonically 
or semi-leptonically. For their production, we have focused on rare decays of 
mesons and charged leptons which are copiously produced at various facilities such as beam-dump and 
collider experiments. Once produced, these light binos can lead to exotic signatures such as 
displaced vertices (DVs) or missing energy.
We have used searches for these signatures to constrain the RPV couplings associated with a light bino.


These various strategies are experimentally widely utilized to constrain heavy neutral leptons 
(HNLs) which may decay to almost the same final-state particles as the lightest neutralino in the 
RPV models. We have thus used the existing HNL searches to set new strict bounds on the relevant 
RPV couplings.
Furthermore we have translated the prjected sensitivity to the HNL parameters at certain future experiments into the corresponding search sensitivity for the light neutralinos.





We have studied comprehensively the past experiments \texttt{PIENU}, 
\texttt{NA62}, \texttt{T2K}, and \texttt{BaBar}, as well as the approved experiments \texttt{FASER}, \texttt{MoEDAL-MAPP1}, \texttt{PIONEER}, and \texttt{DUNE}.
We did not consider future experiments that are not yet approved such as \texttt{MATHUSLA} and \texttt{ANUBIS}, with two exceptions, namely \texttt{FASER2} and \texttt{MoEDAL-MAPP2} since they would be the follow-up programs of the two currently 
running experiments \texttt{FASER} and \texttt{MoEDAL-MAPP1}.


Given the various types and flavor-indices of the RPV operators that can be switched on, we have 
investigated separately different theoretical benchmark scenarios, which can be bounded by distinct 
experiments and strategies. For the selected representative benchmark scenarios, we have performed 
numerical computation and presented the final exclusion bounds. Further, we have compared these 
recast bounds with the existing limits on the RPV couplings which mainly stem from low-energy 
processes of meson and lepton decays. In general, we find that in most cases the exclusion limits 
obtained from recasting past HNL searches surpass the existing bounds on the RPV 
couplings by orders of magnitude, and the expected limits at the considered ongoing and future 
experiments can be even stronger.



Simple and analytic reinterpretation methods are becoming an important research tool.
This is because most published experimental reports present results only for a limited number of simple models, and a recast with full simulation is often complicated and time-consuming.
Some existing works such as Refs.~\cite{Fernandez-Martinez:2023phj, Barouki:2022bkt, Beltran:2023nli} have shown the power of simple and quick reinterpretation of searches for long-lived particles by considering heavy neutral leptons in various models and axion-like particles as examples.
Our work exemplifies again the strengths and convenience of such reinterpretation methods, by recasting the bounds on the HNLs in the minimal scenario, into those on the lightest neutralinos in the RPV supersymmetric models, and hence motivates the development of further studies with these reinterpretation methods.




